\chapter{Recommendations, Conclusion, and Future Works}
\label{chap:recommendations-conclusion-future-works}
This chapter describes our recommendations for the different users of the
system, conclusions as well as future works for the system.
\section{Recommendations}
Although the Traxis system successfully addresses several challenges associated with private shuttle transport coordination, there are opportunities to further enhance its effectiveness and adoption. The following recommendations are proposed:
\begin{itemize}
\item Integration of Mobile Payment Solutions: Future deployments should integrate mobile payment platforms such as MTN Mobile Money, Airtel Money, and banking services to enable passengers to pay for bookings electronically. This would reduce cash handling, minimize booking cancellations, and improve financial accountability for operators.

\item Partnership with Shuttle Operators: To ensure widespread adoption, it is crucial to establish partnerships with private shuttle companies. Operators should be involved in the development and deployment process to ensure the system meets their operational needs and encourages buy-in.

\item Continuous User Training and Awareness: Since users possess varying levels of digital literacy, awareness campaigns and training sessions should be conducted to educate passengers and operators on the benefits and proper use of the application. Simple user guides and customer support channels should also be provided.

\item Improvement of Internet Infrastructure Support: The application relies on internet connectivity for real-time tracking and communication. Developers should incorporate offline synchronization mechanisms and optimize the system for low-bandwidth environments common in many parts of Uganda.

\item Regular System Maintenance and Security Updates: Continuous monitoring, maintenance, and security updates should be implemented to ensure system reliability, protect user data, and address emerging cybersecurity threats.

\item Government and Regulatory Support: Collaboration with transport authorities and relevant government agencies would facilitate wider adoption, improve standardization within the private transport sector, and promote digital transformation initiatives.
\end{itemize}

\section{Conclusion}
This project successfully designed and developed Traxis, a mobile-based coordination system intended to improve the management of private shuttle and courier transport services in Uganda. The study identified that manual booking procedures, poor coordination, limited service visibility, and inefficient seat utilization significantly affect the efficiency and reliability of private shuttle operations.

To address these challenges, the developed system provides structured seat booking, real-time vehicle tracking, passenger scheduling, and operator management through an integrated platform. The dual-interface architecture enables passengers to conveniently access transport services while allowing operators to effectively monitor trips, manage bookings, and optimize vehicle capacity.

The implementation and testing of the system demonstrate that digital technologies can significantly improve operational efficiency, transparency, and customer experience within Uganda's private transport sector. Furthermore, the project establishes a scalable foundation for the modernization of shuttle and courier services through mobile technology.

Overall, the objectives of the study were achieved, and the Traxis system demonstrates the potential to contribute toward a more organized, accessible, and efficient transportation ecosystem.

\section{Future Works}

While Traxis provides a functional solution for private shuttle coordination, several enhancements can be implemented in future versions to increase its capabilities and scalability:
\begin{itemize}

\item Integration of AI-Based Demand Prediction: Machine learning algorithms can be incorporated to predict passenger demand, recommend optimal schedules, and improve vehicle allocation based on historical travel patterns.

\item Expansion to Multiple Routes and Cities: The current implementation focuses primarily on specific intercity routes. Future work should extend the platform to support nationwide operations involving multiple operators and routes across Uganda.

\item Dynamic Route Optimization: Future versions may include intelligent route optimization that considers traffic conditions, passenger requests, and road conditions to minimize travel time and operational costs.

\item Comprehensive Parcel Tracking System: The courier functionality can be expanded to include barcode or QR-code scanning, real-time parcel status updates, estimated delivery times, and proof-of-delivery mechanisms.

\item Integration with Digital Payment Gateways: Support for mobile money, debit/credit cards, and other electronic payment solutions would enable end-to-end digital transactions and automated financial reporting.

\item Advanced Analytics Dashboard: Future systems can provide operators with detailed analytics on revenue, occupancy rates, customer behavior, route performance, and operational trends to support better decision-making.

\item Passenger Rating and Feedback System: A two-way rating mechanism allowing passengers and operators to evaluate each other would improve service quality, accountability, and trust within the platform.

\item Multilingual and Accessibility Support: To increase inclusivity, future versions should support multiple local languages, voice-assisted navigation, and accessibility features for users with disabilities.

\item Integration with Government Transport Information Systems: Future collaboration with regulatory authorities could enable licensing verification, compliance monitoring, and standardized transport information sharing across the sector.

\item Cross-Platform Web Portal: In addition to the mobile application, a comprehensive web-based portal could be developed for passengers and operators to manage bookings, monitor services, and generate reports from desktop environments.
\end{itemize}